\section{Building the intrinsic-alignment model and dataset}
\label{sec:model}

Section~\ref{sec:background} introduced intrinsic alignment as the physical
effect we want to represent.  This section turns that problem into a
controlled dataset for compression.  This construction matters because a
compression score is meaningful only if the input family spans the kinds of
variation we intend to test and remains numerically valid throughout the
sampled region.  We therefore separate the overall amplitude from the shape
of the response, define smooth redshift and scale variation, test the sampling
bounds, and create independent subsets for fitting and evaluation.

\subsection{From three-dimensional IA power to the response amplitude}

Section~\ref{sec:ia-background} introduced GG, GI, and II as contributions to
observed shape correlations.  Their three-dimensional NLA ingredients are the
matter--intrinsic cross-power spectrum and the intrinsic-shape auto-power
spectrum:
\[
  P_{\delta I}(k,z)=\mathcal A_{\rm IA}(k,z)P_\delta^{\rm nl}(k,z),
  \qquad
  P_{II}(k,z)=\mathcal A_{\rm IA}^{\,2}(k,z)P_\delta^{\rm nl}(k,z).
\]
Here $P_\delta^{\rm nl}(k,z)$ is the nonlinear matter power spectrum.  The
first relation connects an intrinsic shape to the surrounding matter field,
whereas the second connects two intrinsic shapes responding to correlated
tidal environments.  The sign of $P_{\delta I}$ follows the signed response
$\mathcal A_{\rm IA}$; $P_{II}$ contains its square.  These spectra are model
ingredients, not quantities read directly from a galaxy catalogue.  Survey
predictions project them along the line of sight with source and lensing
kernels to obtain the angular GI, IG, and II contributions discussed in
Section~\ref{sec:ia-background} \citep{2007NJPh....9..444B}.

\begin{figure}[!htbp]
  \centering
  \includegraphics[width=\linewidth]{ia_power_spectra_panels.pdf}
  \caption{The three power spectra used in the NLA construction at five
  redshifts.  Panel (a) is the nonlinear matter spectrum.  Panel (b) shows
  $-P_{\delta I}$ because the adopted sign convention makes
  $P_{\delta I}<0$; the signed values, rather than their negatives, enter the
  model.  Panel (c) shows the positive $P_{II}$.  The same
  $\mathcal A_{\rm IA}$ enters linearly in panel (b) and quadratically in
  panel (c), which explains their different redshift ordering and amplitude.
  The curves use the notebook's Planck-like parameters and fiducial IA
  response, with an Eisenstein--Hu transfer function and Halofit nonlinear
  correction in PyCCL.  They illustrate the model relations and are not a
  fit to survey data.}
  \label{fig:ia-power-spectra}
\end{figure}

We factor the signed response as

\begin{equation}
  \mathcal A_{\rm IA}(k,z)
  = -A_0\,\mathcal A_\Omega(z)\,\mathcal A_\Theta(k,z).
  \label{eq:ia_decomposition}
\end{equation}

Here $A_0$ sets the overall IA strength.  The standard cosmological factor
$\mathcal A_\Omega(z)$ is proportional to $\Omega_{\rm m}/D(z)$, where
$D(z)$ is the linear growth factor.  The positive surface
$\mathcal A_\Theta(k,z)$ supplies the additional scale and redshift
flexibility tested in this project.  It is defined at the start because it is
the object later passed to PCA and the autoencoder:

\begin{equation}
  \mathcal A_\Theta(k,z)=R_z(z)R_L(z)S(k,z)>0.
  \label{eq:shape_surface}
\end{equation}

An overall multiplier is easy to compress and would say little about whether
a method had learned variation across the surface.  We therefore keep $A_0$
outside the 13-dimensional shape family and compress
$\log_{10}\mathcal A_\Theta$, not the signed full amplitude.  The three
factors in Equation~\ref{eq:shape_surface} have distinct roles: $R_z$
describes a broad redshift trend, $R_L$ adds a smooth transition, and $S$
describes scale dependence.

\subsection{The broad redshift factor \texorpdfstring{$R_z$}{Rz}}

Real lensing surveys contain galaxies observed at different cosmic times, so
an IA response that is fixed with redshift would be too narrow for this test.
We measure evolution relative to the fixed pivot $z_*=0.5$, which gives the
parameters a common reference point.  The factor $R_z$ is a power law in
$1+z$, controlled by $\eta$.  Its purpose is to represent a gradual increase
or decrease across the full redshift interval with one interpretable
parameter.  The exact definition is collected in
Appendix~\ref{app:detailed-equations}.

\subsection{The smooth transition factor \texorpdfstring{$R_L$}{RL}}

A single broad tilt cannot represent a response that changes behaviour around
an intermediate redshift.  The factor $R_L$ adds this second, smooth trend
without assigning an independent amplitude to every redshift bin.  The
parameter $\xi$ sets its strength and direction, $z_q$ locates the transition,
and $s$ controls how sharply it occurs.  The term ``luminosity-like'' describes
the shape of this phenomenological factor; it is not calibrated here to a
specific survey luminosity function.  Its definition is given in
Appendix~\ref{app:detailed-equations}.

\begin{figure}[!htbp]
  \centering
  \includegraphics[width=0.86\linewidth]{ia_redshift_factors_panels.pdf}
  \caption{Redshift-factor parameters varied one at a time.  Panel (a) varies
  the broad trend $\eta$; panel (b) varies $\xi$ at fixed $s=2$ and
  $z_q=1.5$; panel (c) varies $z_q$ at fixed $\xi=1$ and $s=4$.  Every curve
  equals unity at the dotted pivot $z_*=0.5$.  These are phenomenological
  factors, not measured luminosity evolution.}
  \label{fig:ia-redshift-factors}
\end{figure}
\FloatBarrier

\subsection{The joint scale factor \texorpdfstring{$S(k,z)$}{S(k,z)}}

The factors $R_z$ and $R_L$ change only with redshift.  To represent the
additional uncertainty associated with nonlinear physical scales, we also
need a factor that changes across $k$.  Figure~\ref{fig:shape-surface} shows
the resulting two-dimensional $S(k,z)$ for the fiducial parameter choices.
Low $k$ corresponds to large physical scales and remains close to the
baseline $S=1$; the departure appears at higher $k$, where smaller physical
scales are being probed.  Its evolution across the vertical direction allows
the scale response itself to change with cosmic time.

\begin{figure}[!htbp]
  \centering
  \includegraphics[width=0.72\linewidth]{ia_scale_redshift_surface.pdf}
  \caption{Fiducial scale--redshift factor $S(k,z)$.  The horizontal axis is
  $\log_{10}(k/{\rm Mpc}^{-1})$ and the vertical axis is redshift.  At
  $z_*=0.5$, the main choices are $k_{t,*}=0.5\,{\rm Mpc}^{-1}$, $q=1$,
  $n_*=2$, $\alpha=0.3$, and $m=2$.  The following subsections separate the
  role of each parameter and its redshift-evolution index.}
  \label{fig:shape-surface}
\end{figure}
\FloatBarrier

\subsection{The scale-correction strength \texorpdfstring{$q$}{q}}

The parameter $q$ controls the overall size of the scale-dependent
correction: $q=0$ removes that correction and returns $S=1$, while larger
sampled values make it more prominent.  Its operational physical meaning is
therefore the strength of the departure from the large-scale NLA baseline.
Varying $q$ is needed to test whether compression works for both nearly
scale-independent and strongly scale-dependent surfaces.  It does not claim
that the strength of real small-scale IA is known.

\subsection{The transition scale \texorpdfstring{$k_t(z)$}{kt(z)}}

The correction should not affect all wavenumbers equally.  Large scales are
closer to the regime in which the NLA prescription is normally used, whereas
additional uncertainty is expected as $k$ increases.  We therefore use
$k_t(z)$ to mark where the scale correction begins to turn on:
$k\ll k_t$ is close to the large-scale baseline, and $k\gtrsim k_t$ is
increasingly affected by $S$.  The particular smooth transition used here is
a construction of this project, not a new physical law or an observationally
measured boundary.  Its purpose is to represent uncertainty in the onset of
small-scale behaviour, not to solve the underlying small-scale physics.
$k_{t,*}$ fixes the transition at $z_*=0.5$, and $\gamma_t$ allows that
location to move with redshift.

\subsection{The transition sharpness \texorpdfstring{$n(z)$}{n(z)}}

Knowing the location of a transition is not enough: the response may change
gradually across a broad range of scales or relatively rapidly near $k_t$.
The factor controlled by $n(z)$ distinguishes these possibilities.  Smaller
$n$ gives a broader turn-on and larger $n$ gives a sharper one.  Introducing
this freedom avoids imposing one arbitrary transition width on every
synthetic surface.  The pivot value $n_*$ sets the sharpness at $z_*=0.5$,
while $\gamma_n$ allows the width of the transition to evolve.

\subsection{The high-\texorpdfstring{$k$}{k} slope
\texorpdfstring{$\alpha(z)$}{alpha(z)}}

After the correction has turned on, its amplitude need not approach a flat
plateau.  The parameter $\alpha(z)$ therefore controls the asymptotic
high-$k$ trend: positive values allow the response to continue increasing,
negative values allow it to decrease, and $\alpha=0$ gives no additional
high-$k$ tilt.  This factor is included because a compression method should
not succeed only when every small-scale tail has the same slope.  The pivot
value $\alpha$ specifies the slope at $z_*=0.5$, and $\gamma_\alpha$ controls
its redshift evolution.

\subsection{The tail smoothness \texorpdfstring{$m(z)$}{m(z)}}

The high-$k$ slope must join the transition without an unphysical corner.
The parameter $m(z)$ controls how smoothly the model bends toward its
asymptotic tail.  It is kept separate from $n(z)$ because $n$ controls how the
scale correction switches on, whereas $m$ controls the curvature of the tail
after that onset.  This distinction lets the dataset contain surfaces with
similar transition locations but different detailed shapes.  The pivot value
$m$ and evolution index $\gamma_m$ set this behaviour across redshift.  The
complete expression combining $q$, $k_t$, $n$, $\alpha$, and $m$ is given in
Appendix~\ref{app:detailed-equations}.

These quantities are phenomenological shape controls rather than direct
observables.  The pivot parameters specify the scale response at $z_*$, and
the four dimensionless $\gamma$ indices specify how those controls evolve
away from the pivot.  Their physical interpretation is therefore
operational---strength, onset scale, transition width, tail slope, and tail
curvature---rather than a claim that each parameter corresponds to one
independently measurable galaxy property.

\subsection{Sampling and checking the 13-parameter model family}

The 13 parameters comprise four redshift and transition parameters, five
pivot-scale shape parameters, and four parameters controlling redshift
evolution.  A large ensemble is needed because PCA and a neural network should
be tested on the family as a whole, not on a few hand-picked curves.  We draw
100,000 combinations from documented data-generation
bounds; the complete ranges and sampling distributions are listed in
Appendix Table~\ref{tab:sampling-bounds}.  Each combination produces a
surface on 31 redshift points from $z=0$ to $z=3$ and 101 logarithmically
spaced wavenumbers from $k=0.01$ to $10\,\mathrm{Mpc}^{-1}$.  The redshift
grid covers the survey-relevant evolution represented by the model, while
logarithmic spacing resolves changes across several decades in scale.

Random samples alone can miss failures at the edges of a high-dimensional
box.  Before fixing these bounds, we therefore evaluated
8,192 random draws and all $2^{13}=8{,}192$ corners of the parameter box.  All
of those surfaces were finite and positive.  Dynamic-range checks were used
to keep the family broad enough for compression tests without allowing
unmanageably extreme corner cases.

\reportfigure{sampling_amplitude_ensemble.pdf}{Distribution of 8,192 random
sampling-range validation surfaces at $z=0$ and $z=3$.  The shaded regions contain the
central 98\% and 90\% of $\mathcal A_\Theta$ at each $k$, the black curve is
the median, and the colored curves are the five surfaces with the largest
internal dynamic range.  All evaluated surfaces are finite and strictly
positive.}{fig:sampling-ensemble}

\subsection{Training, validation, and test subsets}

Evaluating a model on the same surfaces used to fit it would overstate how
well it generalizes.  The stored split therefore contains 80,000 training
surfaces, 10,000 validation surfaces, and 10,000 test surfaces.  The training
subset fits the PCA basis and autoencoder; the validation subset supports
model selection and diagnostics; and the test subset remains unused until
all choices have been frozen.  Normalization statistics and the PCA basis are
calculated from training data only so that information from held-out surfaces
cannot influence the fitted representations.

This is a phenomenological model family.  The sampled ranges are documented
data-generation bounds, not observational constraints, and the surfaces have
not been calibrated to a particular galaxy survey.  The resulting dataset
therefore gives us a controlled test of compression: every surface has the
same $(k,z)$ grid, but the underlying parameters vary across the prescribed
13-dimensional space.  In the next section, we use this dataset to test
whether the $3{,}131$ grid values can be represented by a much smaller latent
vector.

\clearpage
